\documentclass[12pt,a4paper]{report}

% Packages
\usepackage[utf8]{inputenc}
\usepackage[T1]{fontenc}
\usepackage[lining]{ebgaramond}
\usepackage[cmintegrals,cmbraces]{newtxmath}
\usepackage{microtype}
\usepackage{geometry}
\usepackage{tabularx}
\usepackage{graphicx}
\usepackage{longtable}
\usepackage{booktabs}
\usepackage{array}
\usepackage{listings}
\usepackage{xcolor}
\usepackage{hyperref}
\usepackage{amsmath}
\usepackage{pdflscape}
\usepackage{float}
\usepackage{enumitem}
\usepackage{makecell}
\graphicspath{{./}}

% Page layout
\geometry{
    left=1.8cm,
    right=1.8cm,
    top=2.2cm,
    bottom=2.2cm
}

% Color definitions
\definecolor{crimson}{RGB}{153, 0, 0}
\definecolor{darkblue}{RGB}{0, 0, 153}

% Hyperlink settings
\hypersetup{
    colorlinks=true,
    linkcolor=crimson,
    urlcolor=crimson,
    citecolor=crimson
}

% Allow line breaks at slashes to prevent overfull hbox
\tolerance=2000
\emergencystretch=2em

% Chapter/section title formatting
\usepackage{titlesec}
\titleformat{\chapter}{\normalfont\huge\bfseries}{}{0pt}{}[\vspace{0.3em}]
\titlespacing*{\chapter}{0pt}{0pt}{14pt}
\titleformat{\section}{\normalfont\Large\bfseries}{\thesection}{1em}{}
\titleformat{\subsection}{\normalfont\large\bfseries}{\thesubsection}{1em}{}

\begin{document}

\pagestyle{plain}

% ─────────────────────────────────────────────
%  Cover Page
% ─────────────────────────────────────────────
\begin{titlepage}
\centering
\vspace*{3cm}

{\fontsize{28}{34}\selectfont\bfseries DSP in VLSI\par}
\vspace{0.6em}
{\fontsize{20}{26}\selectfont\bfseries Homework 3\par}
\vspace{0.4em}
{\fontsize{16}{22}\selectfont\bfseries Interpolator\par}

\vspace{1.5cm}
\textcolor{crimson}{\rule{0.6\linewidth}{1.2pt}}
\vspace{1.5cm}

{\large\bfseries Student ID}\par
\vspace{0.4em}
{\large M11407439}\par

\vspace{0.8em}

{\large\bfseries Name}\par
\vspace{0.4em}
{\large MING-HONG, LIN}\par

\vspace{2cm}

{\large\bfseries Development Environment}\par
\vspace{0.8em}
\begin{tabular}{rl}
    \textbf{HDL}              & Verilog HDL \\[4pt]
    \textbf{Simulation Tools} & Synopsys VCS and Verdi \\[4pt]
    \textbf{Synthesis Tool}   & Synopsys DC Compiler \\[4pt]
    \textbf{Process}          & TSMC 16nm(ADFP) \\[4pt]
    \textbf{Verification}     & Python \\
\end{tabular}

\vfill

\textcolor{crimson}{\rule{0.6\linewidth}{0.6pt}}\par
\vspace{0.5em}
{\small Due Date: 2026/04/15}

\end{titlepage}

% ─────────────────────────────────────────────
%  Table of Contents
% ─────────────────────────────────────────────
\tableofcontents
\newpage

% ─────────────────────────────────────────────
\chapter{Step 1 — Interpolator Function Verification (Over-Sampled Cases)}
% ─────────────────────────────────────────────

Assume that
\[
    x_1[m] = \cos\!\left(2\pi\!\left(\frac{mT_s}{10T_s}+\frac{\phi}{2}\right)\right)
            + j\sin\!\left(2\pi\!\left(\frac{mT_s}{10T_s}+\frac{\phi}{2}\right)\right),
\]
where $\phi = \mathrm{mod}(d,\,2)$ and $d$ is the last digit of your student ID.
The sampling frequency is 10 times the frequency of the sinusoidal wave.

Please use \textbf{linear interpolator} and \textbf{second-order polynomial
interpolator} to interpolate the sampled waveform in the region of
$10 \le m \le 20$ with $\mu = 0,\tfrac{1}{8},\tfrac{2}{8},\ldots\tfrac{7}{8}$.
Note that you have $x_1[9]$ when you want to interpolate $x_1[10.125]$.

Show the \textbf{real-part} and \textbf{imaginary-part} output waveform
$\hat{x}_1[m+\mu]$ after interpolation for both methods.
In addition, draw the absolute error
$|x_1[m+\mu]-\hat{x}_1[m+\mu]|$
in the region of $10 \le m \le 20$ with
$\mu = 0,\tfrac{1}{8},\tfrac{2}{8},\ldots\tfrac{7}{8}$.
\textbf{(30\%)}

\section{Linear vs Second-Order Polynomial Interpolator}
\begin{figure}[H]
    \centering
    \includegraphics[width=1\linewidth]{image0.png}
    \caption{Waveform of Linear vs Second-Order Polynomial Interpolator}
    \label{fig:placeholder}
\end{figure}

\begin{figure}[H]
    \centering
    \includegraphics[width=1\linewidth]{image1.png}
    \caption{Absolute Error Comparison of Linear vs Second-Order Polynomial Interpolator}
    \label{fig:placeholder}
\end{figure}

\newpage

% ─────────────────────────────────────────────
\chapter{Step 2 — Interpolator Function Verification (Typical Cases)}
% ─────────────────────────────────────────────

Assume that
\[
    x_2[m] = 2^{-20}\!\left[\cos\!\left(2\pi\!\left(\frac{mT_s}{4T_s}+\frac{\phi}{3}\right)\right)
            + j\sin\!\left(2\pi\!\left(\frac{mT_s}{4T_s}+\frac{\phi}{3}\right)\right)\right],
\]
where $\phi$ is the same parameter according to your student ID.
The sampling frequency is 4 times the frequency of the sinusoidal wave.

Please use \textbf{linear interpolator} and \textbf{piecewise parabolic interpolator}
to interpolate the sampled waveform in the region of $4 \le m \le 8$ with
$\mu = 0,\tfrac{1}{9},\tfrac{2}{9},\ldots\tfrac{8}{9}$.

Show the \textbf{real-part} and \textbf{imaginary-part} output waveform
$\hat{x}_2[m+\mu]$ after interpolation.
In addition, draw the absolute error
$|x_2[m+\mu]-\hat{x}_2[m+\mu]|$
in the region of $4 \le m \le 8$ with
$\mu = 0,\tfrac{1}{8},\tfrac{2}{8},\ldots\tfrac{7}{8}$.
\textbf{(30\%)}

\section{Linear vs Piecewise Parabolic Interpolator}

\begin{figure}[H]
    \centering
    \includegraphics[width=1\linewidth]{image2.png}
    \caption{Waveform of Linear vs Piecewise Parabolic Interpolator}
    \label{fig:placeholder}
\end{figure}

\begin{figure}[H]
    \centering
    \includegraphics[width=1\linewidth]{image3.png}
    \caption{Absolute Error Comparison of Linear vs Piecewise Parabolic Interpolator}
    \label{fig:placeholder}
\end{figure}
\newpage

% ─────────────────────────────────────────────
\chapter{Step 3 — Bit-True Model (BF16)}
% ─────────────────────────────────────────────

Please use \textbf{BF16 format} to realize addition and multiplication.
Write the bit-true model (simulated program in Python) to realize BF16 additions
and multiplication. For simplicity, NaN, infinity, and subnormal numbers are not
handled — if the magnitude is smaller than the smallest normal number, use
\textbf{truncation} to make it become 0. Truncation is also used to avoid
renormalization.

Show the results calculated by your BF16 bit-true model for the following operands.
Express results both in \textbf{decimal} and \textbf{binary} representation
(S = sign bit, E = exponent field, F = fraction field, all in binary).
\textbf{(20\%)}

\section{BF16 Operation Results}

\begin{table}[H]
    \centering
    \renewcommand{\arraystretch}{1.3}
    \begin{tabular}{cllc}
        \toprule
        & \textbf{Operand 1} & \textbf{Operand 2} & \textbf{Operation} \\
        \midrule
        (a) &
        \makecell[l]{S: 1 \\ E: \texttt{1000\_0011} \\ F: \texttt{0000\_000}} &
        \makecell[l]{S: 0 \\ E: \texttt{1000\_1011} \\ F: \texttt{0000\_011}} &
        Addition \\
        \midrule
        (b) &
        \makecell[l]{S: 1 \\ E: \texttt{0000\_0001} \\ F: \texttt{0000\_011}} &
        \makecell[l]{S: 0 \\ E: \texttt{0000\_0001} \\ F: \texttt{1111\_010}} &
        Addition \\
        \midrule
        (c) &
        \makecell[l]{S: 1 \\ E: \texttt{0000\_0010} \\ F: \texttt{1100\_000}} &
        \makecell[l]{S: 0 \\ E: \texttt{0111\_1100} \\ F: \texttt{0000\_110}} &
        Multiplication \\
        \midrule
        (d) &
        \makecell[l]{S: 0 \\ E: \texttt{0110\_0011} \\ F: \texttt{1011\_110}} &
        \makecell[l]{S: 0 \\ E: \texttt{1001\_0011} \\ F: \texttt{0101\_000}} &
        Multiplication \\
        \bottomrule
    \end{tabular}
    \caption{BF16 operands and operations.}
    \label{tab:bf16-ops}
\end{table}

\begin{table}[H]
    \centering
    \renewcommand{\arraystretch}{1.3}
    \begin{tabular}{clllc}
        \toprule
        & \textbf{Operation} & \textbf{Decimal} & \textbf{S} & \textbf{E \& F (binary)} \\
        \midrule
        (a) & Addition       & $4192$      & 0 & E: \texttt{1000\_1011},\ F: \texttt{0000\_011} \\
        (b) & Addition       & $0$         & 0 & E: \texttt{0000\_0000},\ F: \texttt{0000\_000} \\
        (c) & Multiplication & $0$         & 0 & E: \texttt{0000\_0000},\ F: \texttt{0000\_000} \\
        (d) & Multiplication & $0.0088501$ & 0 & E: \texttt{0111\_1000},\ F: \texttt{0010\_001} \\
        \bottomrule
    \end{tabular}
    \caption{BF16 operation results.}
    \label{tab:bf16-results}
\end{table}

\newpage

% ─────────────────────────────────────────────
\chapter{Step 4 — Performance Degradation by Finite Precision}
% ─────────────────────────────────────────────

Please analyze the error of the \textbf{second-order interpolator} if its arithmetic
operations are changed to BF16 by using the bit-true model.

Draw the figures indicating the difference of interpolated results caused by data
format for outputs $\hat{x}_1[m+\mu]$ in the region of $10 \le m \le 20$
with $\mu = 0,\tfrac{1}{8},\tfrac{2}{8},\ldots\tfrac{7}{8}$.
\textbf{(10\%)}

\section{Quantization Error Analysis}
As shown in the figure below, switching the arithmetic from double precision to
BF16 introduces a noticeable quantisation error due to the reduced mantissa
width (7 bits) of the BF16 format.  The peak absolute errors observed over the
interpolation range $10 \le m \le 20$ are summarised below:

\begin{table}[H]
    \centering
    \renewcommand{\arraystretch}{1.3}
    \begin{tabular}{lc}
        \toprule
        \textbf{Component} & \textbf{Max Absolute Error} \\
        \midrule
        Real part & $1.5959 \times 10^{-2}$ \\
        Imaginary part & $1.6048 \times 10^{-2}$ \\
        \bottomrule
    \end{tabular}
    \caption{Maximum absolute error introduced by BF16 arithmetic in the second-order interpolator.}
    \label{tab:bf16-error}
\end{table}

\begin{figure}[H]
    \centering
    \includegraphics[width=1\linewidth]{image4.png}
    \caption{Interpolation Error Caused by BF16 Finite Precision (Second-Order Interpolator)}
    \label{fig:placeholder}
\end{figure}

\newpage

% ─────────────────────────────────────────────
\chapter{Step 5 — Hardware Reduction}
% ─────────────────────────────────────────────

Draw the complete block diagram of your first version implementation according to Fig. 4. (5\%) Show your hardware sharing version (The precision of the datapath must be BF16 and can not be changed.) (10\%) Draw the block diagram after your improvement. (5\%) Indicate the critical path (5\%) 

\section{Original Block Diagram (Farrow Structure)}

The second-order polynomial interpolator in Farrow structure decomposes the
interpolation into three fixed sub-filters $v_0$, $v_1$, and $v_2$, whose
outputs are combined as
\[
    \hat{x}[m+\mu] = v_0 + \mu\,v_1 + \mu^2\,v_2.
\]
The coefficients of each sub-filter are derived from three consecutive input
samples $X[m]$, $X[m+1]$, and $X[m+2]$ as summarised in
Table~\ref{tab:farrow-coeffs}.

\begin{table}[H]
    \centering
    \renewcommand{\arraystretch}{1.4}
    \begin{tabular}{clll}
        \toprule
        \textbf{Sub-filter} & \textbf{$X[m]$} & \textbf{$X[m+1]$} & \textbf{$X[m+2]$} \\
        \midrule
        $v_2$ & $+\tfrac{1}{2}$ & $-1$   & $+\tfrac{1}{2}$ \\
        $v_1$ & $-\tfrac{3}{2}$ & $+2$   & $-\tfrac{1}{2}$ \\
        $v_0$ & $+1$            & $\;0$  & $\;0$           \\
        \bottomrule
    \end{tabular}
    \caption{Farrow sub-filter coefficients for the second-order polynomial interpolator.}
    \label{tab:farrow-coeffs}
\end{table}

\begin{figure}[H]
    \centering
    \includegraphics[width=1\linewidth]{Interpolator.png}
    \caption{Enter Caption}
    \label{fig:placeholder}
\end{figure}

\newpage

\section{Hardware Sharing Approaches}

Three hardware optimisations were applied to reduce resource usage and improve
operating frequency.

\begin{enumerate}[leftmargin=*, label=\textbf{\arabic*.}]

\item \textbf{Shared delay-line (single DFF chain).}
      The original structure uses three independent delay lines, one per
      sub-filter branch.  Since all three branches consume the same input
      samples $X[m]$, $X[m+1]$, and $X[m+2]$, a single shared shift-register
      chain suffices, eliminating two redundant sets of D flip-flops.

\item \textbf{Multiplier-free scaling.}
      The dominant coefficient in the design is $0.5$.  In BF16
      format, multiplication by $0.5$ reduces to decrementing the
      exponent field by one — no multiplier is needed.  Accordingly, all
      $\times 0.5$ operations are replaced by a one-bit exponent
      adjustment.  The coefficient $\times 2$ is likewise converted to an
      equivalent addition ($A + A$), and $\times 1.5$ is decomposed
      into $A + 0.5A$, so that the entire datapath requires no
      general-purpose multipliers.

\item \textbf{Two-stage pipeline insertion.}
      To increase the maximum operating frequency, two pipeline registers are
      inserted across the critical path, balancing the combinational logic
      depth between stages.

\end{enumerate}

\begin{figure}[H]
    \centering
    \includegraphics[width=1\linewidth]{Interpolator2.png}
    \caption{Block Diagram with Hardware Sharing Approaches}
    \label{fig:placeholder}
\end{figure}

\newpage

% ─────────────────────────────────────────────
\chapter{Step 6 — Implementation}
% ─────────────────────────────────────────────

Please write Verilog to implement the \textbf{second-order interpolator} of
Farrow structure using \textbf{BF16 arithmetic units} after hardware reduction.
Insert D flip-flops at inputs and outputs.

Note: The input to be interpolated changes \textbf{every 8 clock cycles};
the $\mu$ value changes \textbf{every clock cycle}.

\section{Behavior Simulation}

Show the timing diagram of behavior simulation for the
\textbf{first 8} ($\mu = 0,\tfrac{1}{8},\ldots\tfrac{7}{8}$) and
\textbf{last 8} ($\mu = 0,\tfrac{1}{8},\ldots\tfrac{7}{8}$)
interpolated outputs of $\hat{x}_1[m+\mu]$
in the region of $10 \le m \le 20$.

Show the \textbf{real-part} and \textbf{imaginary-part} error of hardware
implementation results compared to the double-precision results.
\textbf{(30\%)}

\newpage

\section{Timing Report and Constraint}

Show your \textbf{timing report} to verify the critical path of the block diagram
in Step~5. \textbf{(10\%)}
List your \textbf{timing constraint} for post-synthesis simulation. \textbf{(5\%)}

% [Insert timing report screenshot here]

\begin{table}[H]
    \centering
    \begin{tabular}{lc}
        \toprule
        \textbf{Parameter} & \textbf{Value (ns)} \\
        \midrule
        Clock Period                  & 5\\
        Clock Uncertainty             & 0.1\\
        Clock Latency                 & 0.1\\
        Input / Output Maximum Delay  & 0.5\\
        \bottomrule
    \end{tabular}
    \caption{Timing constraint settings.}
    \label{tab:timing-constraint}
\end{table}

\newpage

\section{Post-Synthesis Simulation}

Show the timing diagram of \textbf{post-synthesis simulation} for the
first 8 and last 8 interpolated outputs of $\hat{x}_1[m+\mu]$
in the region of $10 \le m \le 20$.

Show the \textbf{real-part} and \textbf{imaginary-part} error of post-synthesis
hardware outputs compared to the double-precision results. \textbf{(20\%)}

% [Insert post-synthesis timing diagram (first 8 and last 8 outputs) here]
% [Insert real-part and imaginary-part error figures here]

\newpage

\section{BF16 Bit-True Model vs. Verilog Output Verification}

Construct the \textbf{BF16 bit-true model} following the architecture after
hardware improvement. Depict the \textbf{real-part} and \textbf{imaginary-part}
errors of interpolated outputs $\hat{x}_1[m+\mu]$ between the Verilog outputs
and the bit-true model in the region of $10 \le m \le 20$
with $\mu = 0,\tfrac{1}{8},\tfrac{2}{8},\ldots\tfrac{7}{8}$. \textbf{(20\%)}
\begin{figure}[H]
    \centering
    \includegraphics[width=1\linewidth]{image5.png}
    \caption{BF16 Bit-True Model vs. Verilog Output}
    \label{fig:placeholder}
\end{figure}

\end{document}
